\section*{Chapitre \ref{ch:dv:the-structured-products-business} --- Le métier des produits structurés}

\begin{solution}{exo:dv:the-structured-products-business:1}
Une obligation zéro-coupon de l'émetteur, versant 100 dans deux ans, et 0.6 calls à la monnaie sur 100 de l'indice venant à échéance dans deux ans.
L'investisseur détient les deux~: il est long de l'obligation de l'émetteur (donc de son risque de crédit) et long des calls.
\end{solution}

\begin{solution}{exo:dv:the-structured-products-business:2}
$100e^{-0.06}=94.18$ à spread nul et $100e^{-0.108}=89.76$ à 2.4~\%~: le spread ajoute 4.41 de budget d'option, plus que doublant
les 4.32 restant à spread nul après la marge.
\end{solution}

\begin{solution}{exo:dv:the-structured-products-business:3}
Une participation plus élevée peut être payée par le spread de financement plus élevé d'un émetteur moins solide, qui rémunère l'investisseur d'un risque
plus élevé de perdre la protection en cas de défaut~; ou par une \omterm{def:dv:the-structured-products-business:margin}{marge de structuration} plus large prélevée ailleurs. Il faut comparer les notes au même crédit d'émetteur, ou ajouter la
valeur du risque de crédit à la comparaison.
\end{solution}

\begin{solution}{exo:dv:the-structured-products-business:4}
Le budget est $100-0.95\times94.18-1.5=9.03$, la participation vaut donc $9.03/12.49=72.3\%$~: cinq points de protection achètent plus
que le double de participation.
\end{solution}

\begin{solution}{exo:dv:the-structured-products-business:5}
La volatilité de l'indice est maintenue près de la cible par construction, si bien que l'option est évaluée près de la volatilité cible, bien en dessous de
celle du sous-jacent, et son \omterm{def:dv:greeks-and-the-hedging-pnl:greeks}{vega} est faible. L'investisseur renonce à l'exposition quand la volatilité est élevée (après des baisses, typiquement) et subit
le rééquilibrage de la règle, qui vend après les baisses et achète après les hausses~; les rendements dans les marchés tendanciels diffèrent de ceux de l'indice ordinaire.
\end{solution}

\begin{solution}{exo:dv:the-structured-products-business:6}
Il accuse un retard égal à l'excédent du décrément sur les dividendes, soit 2~\% par an. Si les dividendes tombent à 1~\%, l'indice de cours gagne par rapport
à l'indice de rendement total, alors que l'\omterm{def:dv:the-structured-products-business:decrement}{indice à décrément} ne modifie pas sa règle~: le retard s'élargit à 4~\% par an, supporté par l'investisseur dans l'\omterm{def:dv:the-structured-products-business:decrement}{indice à décrément}~;
l'émetteur, couvert sur le décrément fixe, n'est pas affecté.
\end{solution}

\begin{solution}{exo:dv:the-structured-products-business:7}
Avec un spread de 1~\% et une marge de 1.5~: 5.17~\% avec un strike à 80~\%, 6.96~\% à 90~\% et 9.71~\% à 100~\%. Le bénéfice de financement est le même, 3.92~;
les puts vendus valent 2.55, 4.26 et 6.91 pour 100.
\end{solution}

\begin{solution}{exo:dv:the-structured-products-business:8}
La participation dépend au moins autant du spread de financement de l'émetteur (et de la marge) que des prix des options~: à prix d'options identiques, un émetteur fragile
offre davantage de participation. L'affirmation confond le crédit et la valorisation.
\end{solution}

\begin{solution}{pb:dv:the-structured-products-business:1}
\textbf{1.} 20.5~\%~; 12.49 pour 100.
\textbf{2.} $0.7\times12.49=8.74$.
\textbf{3.} Le call est fixé au niveau initial, à la monnaie~; son prix utilise la volatilité du smile à ce strike et à cette échéance.
\textbf{4.} Le réduire, à peu près dans la mesure où la volatilité du niveau moyenné est plus faible (chapitre 16), ce qui relève la participation.
\textbf{5.} Revendre la hausse au-dessus du cap, ce qui abaisse le prix de l'option et relève la participation en dessous du cap.
\textbf{6.} 94.18.
\textbf{7.} 34.6~\%, 49.6~\% et 64.2~\%.
\textbf{8.} 2.40~\% par an (240 points de base)~; l'obligation coûte alors 89.76.
\textbf{9.} La \omterm{def:dv:the-structured-products-business:note}{note structurée} est la dette de l'émetteur~: les 100 garantis sont une promesse de l'émetteur, valorisée au taux auquel celui-ci emprunte.
\textbf{10.} Le risque de crédit de l'émetteur~: en cas de défaut, la protection et la valeur des options sont perdues avec l'obligation.
\textbf{11.} 113.3~\%~: participation de 100~\% jusqu'à un gain de 13.3~\%.
\textbf{12.} Un coupon de 6.96~\% par an, le capital étant exposé sous 90~\%.
\textbf{13.} Environ 62~\% à spread nul~: les calls à la monnaie à la volatilité cible de 10~\% coûtent environ 6.94 pour 100 sur deux ans.
\textbf{14.} Elle rend la note généreuse en apparence (participation élevée) et supprime, pour l'émetteur, le risque de volatilité et le \omterm{def:dv:dividends-borrow-and-forwards:risk}{risque de dividende}~; l'investisseur subit la
règle de l'indice, et des rendements qui peuvent différer largement de ceux de l'indice ordinaire.
\textbf{15.} Dans l'Union européenne~: un document d'informations clés d'au plus trois pages A4, autonome, exact, équitable, clair et non
trompeur, avec des scénarios et des coûts standardisés.
\textbf{16.} Il est équitable si 240 points de base correspondent au spread de crédit réel de l'émetteur et si le client souhaite ce risque de crédit~; ce n'est pas un cadeau.
\textbf{17.} Que l'écart correspond au crédit de l'émetteur, et qu'il faut comparer la \omterm{def:dv:the-structured-products-business:note}{note structurée} aux obligations de l'émetteur.
\textbf{18.} Elle la réduit~: la composante obligataire est réévaluée au spread élargi, même si l'indice n'a pas bougé.
\textbf{19.} 2.40~\% par an (240 points de base)~: l'obligation zéro-coupon coûte alors 89.76, laissant 8.74 pour 70~\% du call de 12.49 après la marge de 1.5.
\textbf{20.} Quelle quantité d'option le financement de l'émetteur et la marge laissent à l'investisseur, par unité de hausse du sous-jacent.
\end{solution}

\begin{solution}{iq:dv:the-structured-products-business:1}
$100=\pi\,100e^{-(r+s)T}+p\,C+m$~: le montant protégé est l'obligation zéro-coupon de l'émetteur, la participation multipliée par le prix de la composante optionnelle, et
la marge. La participation augmente avec les taux, avec le spread de l'émetteur et avec une protection moindre, et baisse avec la volatilité, les dividendes
(par le call) et la marge.
\iqlookfor{l'identité et le sens de variation de chaque donnée.}
\end{solution}

\begin{solution}{iq:dv:the-structured-products-business:2}
La banque reçoit aujourd'hui le prix d'émission et doit le paiement plus tard~; couvrir ce paiement avec des dérivés laisse un prêt au
spread de financement propre de la banque. Les \omterm{def:dv:the-structured-products-business:note}{notes structurées} diversifient les sources de financement de la banque et sont valorisées à partir de sa courbe de financement.
\iqlookfor{la \omterm{def:dv:the-structured-products-business:note}{note structurée} comme dette, et le résiduel couvert.}
\end{solution}

\begin{solution}{iq:dv:the-structured-products-business:3}
Une \omterm{def:dv:the-structured-products-business:note}{note structurée} versant un coupon amélioré et remboursant la valeur du sous-jacent, plutôt que le pair, si elle finit sous un strike~: l'investisseur est vendeur d'un put
(et acheteur de l'obligation de l'émetteur). Le coupon est égal au taux de financement plus la prime du put, moins la marge.
\iqlookfor{le put vendu et l'origine du coupon.}
\end{solution}

\begin{solution}{iq:dv:the-structured-products-business:4}
La volatilité de l'indice est pilotée à la cible, si bien que les options sur celui-ci sont valorisées près de la cible et ont peu de \omterm{def:dv:greeks-and-the-hedging-pnl:greeks}{vega}. Ce qui reste pour l'émetteur
est l'écart entre la volatilité estimée par la règle et la volatilité réalisée (sauts brusques avant que l'exposition ne s'ajuste) et les coûts de rééquilibrage de l'indice~;
l'émetteur supporte aussi le risque de modèle lié à la dépendance au chemin de l'indice.
\iqlookfor{la volatilité pilotée et les risques résiduels.}
\end{solution}

\begin{solution}{iq:dv:the-structured-products-business:5}
La composante obligataire de chaque note en circulation vaut moins (actualisée au spread élargi), si bien que les cours secondaires baissent~; un émetteur qui porte ses notes à
fair value constate un gain sur sa propre dette. Les nouvelles émissions peuvent offrir de meilleures conditions. Les plans de financement changent si les rachats augmentent.
\iqlookfor{cours secondaires, comptabilisation du risque de crédit propre et nouvelles émissions.}
\end{solution}

\begin{solution}{iq:dv:the-structured-products-business:6}
Vendre davantage d'optionnalité~: un put plus proche de la monnaie (5.2~\%, 7.0~\%, 9.7~\% pour des strikes de 80~\%, 90~\%, 100~\% dans le chapitre), une barrière quotidienne, un worst-of sur
plusieurs sous-jacents (corrélation), un autocall qui raccourcit la note~; faire appel à un émetteur plus fragile (crédit)~; réduire la marge. Chacun relève le coupon en vendant
davantage de risque à l'investisseur~; il faut expliquer chacun en ces termes.
\iqlookfor{les leviers et le risque transféré par chacun.}
\end{solution}
